\chapter*{บรรณานุกรม}
\addcontentsline{toc}{chapter}{บรรณานุกรม}

\begin{hangparas}{1.25cm}{1}
จิรภัทร จันทมาลี. (2569). \textit{จุลชีววิทยาสำหรับการเกษตร: นิเวศวิทยาของดินและการควบคุมโดยชีววิธี}. จันทบุรี: คณะเทคโนโลยีการเกษตร มหาวิทยาลัยราชภัฏรำไพพรรณี.

จิรภัทร จันทมาลี, และ ชีวะ ทัศนา. (2568). การประยุกต์ใช้ปัญญาประดิษฐ์ร่วมกับเมตาเจโนมิกส์ในการประเมินดัชนีสุขภาพดินสวนทุเรียนภาคตะวันออก. \textit{วารสารวิทยาศาสตร์และเทคโนโลยี มหาวิทยาลัยราชภัฏรำไพพรรณี}, 12(2), 45--58.

ชีวะ ทัศนา. (2569). \textit{ฟิสิกส์แผนใหม่: การจำลองเสมือนจริง AR/XR ร่วมกับการคำนวณเชิงตัวเลขด้วย Python}. จันทบุรี: มหาวิทยาลัยราชภัฏรำไพพรรณี.

กรมวิชาการเกษตร. (2565). \textit{คู่มือการใช้จุลินทรีย์ปฏิปักษ์ไตรโคเดอร์มาเพื่อควบคุมโรคพืชในไม้ผลเขตร้อน}. กรุงเทพฯ: กระทรวงเกษตรและสหกรณ์.

Atlas, R. M., \& Bartha, R. (1998). \textit{Microbial Ecology: Fundamentals and Applications} (4th ed.). Menlo Park, CA: Benjamin/Cummings.

Beijerinck, M. W. (1888). Die Bacterien der Papilionaceen-Knöllchen. \textit{Botanische Zeitung}, 46, 725--735.

Compant, S., Samad, A., Faist, H., \& Sessitsch, A. (2019). A review on the plant microbiome: Ecology, functions, and emerging trends in microbial application. \textit{Journal of Advanced Research}, 19, 29--37.

Glick, B. R. (2014). Bacteria with ACC deaminase can promote plant growth and help to feed the world. \textit{Microbiological Research}, 169(1), 30--39.

Harman, G. E., Howell, C. R., Viterbo, A., Chet, I., \& Lorito, M. (2004). \textit{Trichoderma} species—opportunistic, avirulent plant symbionts. \textit{Nature Reviews Microbiology}, 2(1), 43--56.

Hiltner, L. (1904). Über neuere Erfahrungen und Probleme auf dem Gebiete der Bodenbakteriologie. \textit{Arbeiten der Deutschen Landwirtschaftlichen Gesellschaft}, 98, 59--78.

Lugtenberg, B., \& Kamilova, F. (2009). Plant-growth-promoting rhizobacteria. \textit{Annual Review of Microbiology}, 63, 541--556.

Paul, E. A. (2014). \textit{Soil Microbiology, Ecology and Biochemistry} (4th ed.). Academic Press.

Smith, S. E., \& Read, D. J. (2008). \textit{Mycorrhizal Symbiosis} (3rd ed.). London: Academic Press.

van der Heijden, M. G. A., Bardgett, R. D., \& van Straalen, N. M. (2008). The unseen majority: Soil microbes as drivers of plant diversity and productivity in terrestrial ecosystems. \textit{Ecology Letters}, 11(3), 296--310.

Waksman, S. A. (1952). \textit{Soil Microbiology}. New York: John Wiley \& Sons.

Winogradsky, S. (1890). Recherches sur les organismes de la nitrification. \textit{Annales de l'Institut Pasteur}, 4, 213--231.
\end{hangparas}

\clearpage
